\input{../tex-inputs/latex-korrekturansicht-vorspann}

\section[Berta Zuckerkandl an Arthur Schnitzler, 25. 7. {[}1911{]}]{L04009 Berta Zuckerkandl an Arthur Schnitzler, 25. 7. {[}1911{]}}
\nopagebreak\mylabel{L04009v}
\rehead{ }\normalsize\beginnumbering\briefempfaengerindex{Schnitzler, Arthur@\textsc{Schnitzler, Arthur}!zzzZuckerkandl, Berta@\emph{von Berta Zuckerkandl}!1911-07-251@{25. 7. {[}1911{]}}|(be}
\toendnotes[C]{\smallbreak\pagebreak[2]}
\correspDesc{Versand  durch Berta Zuckerkandl am 25. 7. [1911] in Axenstein
\newline{}Erhalt  durch Arthur Schnitzler am [1. 8. 1911?] in Wien}\toendnotes[C]{\smallbreak}
\Standort{CUL, Schnitzler, B 200.}
\physDesc{Brief, 2 Blätter, 8 Seiten, 2043 Zeichen
\newline{}Handschrift: schwarze Tinte, lateinische Kurrent (\noindent{}Nummerierung des zweiten Bogens: »II.«)
\newline{}Schnitzler: mit Bleistift Vermerk: »\textsc{Zuckerkand}\textcolor{gray}{l}« }\toendnotes[C]{\smallbreak}
\pstart
           \raggedleft{}{\pb}\textcolor{gray}{\textbf{\textcolor{pink}{}\oindex{Hotel Axenstein@\textbf{Hotel Axenstein}, \emph{Hotel}|pw}{}\ledrightnote{\textcolor{pink}{Hotel Axenstein}}PARK HOTEL AXENSTEIN}}\pend
           
\pstart
           \raggedleft{}\textcolor{gray}{\textbf{\textcolor{pink}{VIERWALDSTÄTTERSEE}\oindex{Vierwaldstättersee@\textbf{Vierwaldstättersee}, \emph{See}|pw}{}\ledrightnote{\textcolor{pink}{Vierwaldstättersee}}}}\pend
           
\pstart
           \raggedleft{}\textcolor{gray}{\textbf{\textcolor{pink}{Axenstein}\oindex{Axenstein@\textbf{Axenstein}, \emph{Ausflugsziel}|pw}{}\ledrightnote{\textcolor{pink}{Axenstein}},}}{ }25\textsuperscript{ter} Juli\pend
           
\pstart{}Sehr verehrter Herr Doktor!\pend\vspace{0.5em}
\pstart
           Als ich von \textcolor{pink}{Paris}\oindex{Paris@\textbf{Paris}, \emph{Hauptstadt}|pw}{}\ledrightnote{\textcolor{pink}{Paris}} gegen 8\textsuperscript{ten} Juli wegfuhr stand die \textcolor{green}{Medardus}\pwindex{Schnitzler, Arthur 15. 5. 1862 Wien – 21. 10. 1931 ebd.@\textsc{Schnitzler, Arthur} (15. 5. 1862 Wien – 21. 10. 1931 ebd.), \emph{Schriftsteller, Mediziner}!junge Medardus. Dramatische Historie in einem Vorspiel und fünf Aufzügen@\strich\emph{Der junge Medardus. Dramatische Historie in einem Vorspiel und fünf Aufzügen}|pw}{}\ledrightnote{\textcolor{green}{Der junge Medardus. Dramatische Historie in einem Vorspiel und fünf Aufzügen}}-Angelegenheit so gut – dass ich schon an Sieg
               glaubte – und Ihnen ein Telegram senden wollte. Es war besser ich {\pb}tat es nicht – denn so ist Ihnen
               wenigstens die Enttäuschung erspart worden – welche mir zu Teil wurde. Denn
                  gestern erhielt ich den Bescheid des Direktor \textcolor{blue}{Herz}\pwindex{Hertz, Henri 17.\,6.\,1875 Nogent-sur-Seine – 11.\,10.\,1966 16. arrondissement [Paris]@\textsc{Hertz, Henri} (17.\,6.\,1875 Nogent-sur-Seine – 11.\,10.\,1966 16. arrondissement [Paris]), \emph{Schriftsteller, Journalist, Theaterdirektor}|pw}{}\ledrightnote{\textcolor{blue}{Henri Hertz}} hieher, dass er trotzdem sie \label{K_L04009-1v}\edtext{\begin{otherlanguage}{french}entusiasmés de la \textcolor{green}{pièce}\pwindex{Schnitzler, Arthur 15. 5. 1862 Wien – 21. 10. 1931 ebd.@\textsc{Schnitzler, Arthur} (15. 5. 1862 Wien – 21. 10. 1931 ebd.), \emph{Schriftsteller, Mediziner}!junge Medardus. Dramatische Historie in einem Vorspiel und fünf Aufzügen@\strich\emph{Der junge Medardus. Dramatische Historie in einem Vorspiel und fünf Aufzügen}|pwv}{}\ledrightnote{{$\rightarrow$}\emph{\textcolor{green}{Der junge Medardus. Dramatische Historie in einem Vorspiel und fünf Aufzügen}}}\end{otherlanguage}}{\lemma{\textnormal{\emph{entusiasmés de la pièce}}}\Cendnote{\textnormal{\begin{otherlanguage}{french}enthousiasmés de la pièce\end{otherlanguage}, französisch: begeistert vom Stück}}}\label{K_L04009-1} wären, dies\substVorne{}\textsuperscript{es}\substDazwischen{}e\substHinten{} nach reiflicher Erwägung {\pb}nicht
               zur Aufführung bringen könnten. Und zwar weil I\textsuperscript{stns} Für
               diese Saison 1911–12 die \textcolor{brown}{Porte St.
                  Martin}\orgindex{Théâtre de la Porte Saint-Martin@Théâtre de la Porte Saint-Martin|pw}{}\ledrightnote{\textcolor{brown}{Théâtre de la Porte Saint-Martin}} kontraktlich bereits abgeschlossen habe: Ein \textcolor{green}{Lustspiel}\pwindex{Flers, Robert de 25.\,11.\,1872 Pont-l'Évêque – 30.\,7.\,1927 Vittel@\textsc{Flers, Robert de} (25.\,11.\,1872 Pont-l'Évêque – 30.\,7.\,1927 Vittel), \emph{Schriftsteller}!Primerose, comédie en 3 actes@\strich\emph{Primerose, comédie en 3 actes}|pwv}\pwindex{Caillavet, Gaston Arman de 13.\,3.\,1869 Paris – 13.\,1.\,1915 Saint-Médard-d'Excideuil@\textsc{Caillavet, Gaston Arman de} (13.\,3.\,1869 Paris – 13.\,1.\,1915 Saint-Médard-d'Excideuil), \emph{Schriftsteller}!Primerose, comédie en 3 actes@\strich\emph{Primerose, comédie en 3 actes}|pwv}{}\ledrightnote{{$\rightarrow$}\emph{\textcolor{green}{Primerose, comédie en 3 actes}}} von \textcolor{blue}{Flers}\pwindex{Flers, Robert de 25.\,11.\,1872 Pont-l'Évêque – 30.\,7.\,1927 Vittel@\textsc{Flers, Robert de} (25.\,11.\,1872 Pont-l'Évêque – 30.\,7.\,1927 Vittel), \emph{Schriftsteller}|pw}{}\ledrightnote{\textcolor{blue}{Robert de Flers}} u \textcolor{blue}{Caillavet}\pwindex{Caillavet, Gaston Arman de 13.\,3.\,1869 Paris – 13.\,1.\,1915 Saint-Médard-d'Excideuil@\textsc{Caillavet, Gaston Arman de} (13.\,3.\,1869 Paris – 13.\,1.\,1915 Saint-Médard-d'Excideuil), \emph{Schriftsteller}|pw}{}\ledrightnote{\textcolor{blue}{Gaston Arman de Caillavet}}; \uline{zwei}{ }\label{K_L04009-2v}\edtext{\textcolor{green}{Stücke}\pwindex{Bataille, Henri 4.\,4.\,1872 Nîmes – 2.\,3.\,1922 Rueil-Malmaison@\textsc{Bataille, Henri} (4.\,4.\,1872 Nîmes – 2.\,3.\,1922 Rueil-Malmaison), \emph{Schriftsteller}!femme nue, pièce en quatre actes@\strich\emph{La femme nue, pièce en quatre actes}|pw}\pwindex{Bataille, Henri 4.\,4.\,1872 Nîmes – 2.\,3.\,1922 Rueil-Malmaison@\textsc{Bataille, Henri} (4.\,4.\,1872 Nîmes – 2.\,3.\,1922 Rueil-Malmaison), \emph{Schriftsteller}!Flambeaux, pièce en trois actes@\strich\emph{Les Flambeaux, pièce en trois actes}|pw}{}\ledrightnote{\textcolor{green}{La femme nue, pièce en quatre actes}{\newline}\textcolor{green}{Les Flambeaux, pièce en trois actes}} von \textcolor{blue}{Bataille}\pwindex{Bataille, Henri 4.\,4.\,1872 Nîmes – 2.\,3.\,1922 Rueil-Malmaison@\textsc{Bataille, Henri} (4.\,4.\,1872 Nîmes – 2.\,3.\,1922 Rueil-Malmaison), \emph{Schriftsteller}|pw}{}\ledrightnote{\textcolor{blue}{Henri Bataille}}}{\lemma{\textnormal{\emph{Stücke von Bataille}}}\Cendnote{\textnormal{Am 5. 10. 1911 hatte \textcolor{blue}{Henri Batailles}\pwindex{Bataille, Henri 4.\,4.\,1872 Nîmes – 2.\,3.\,1922 Rueil-Malmaison@\textsc{Bataille, Henri} (4.\,4.\,1872 Nîmes – 2.\,3.\,1922 Rueil-Malmaison), \emph{Schriftsteller}|pwk} Schauspiel \emph{\textcolor{green}{La femme nue}\pwindex{Bataille, Henri 4.\,4.\,1872 Nîmes – 2.\,3.\,1922 Rueil-Malmaison@\textsc{Bataille, Henri} (4.\,4.\,1872 Nîmes – 2.\,3.\,1922 Rueil-Malmaison), \emph{Schriftsteller}!femme nue, pièce en quatre actes@\strich\emph{La femme nue, pièce en quatre actes}|pwk}} am \textcolor{pink}{Théâtre
                     de la Porte Saint-Martin}\oindex{Théâtre de la Porte Saint-Martin@\textbf{Théâtre de la Porte Saint-Martin}, \emph{Theater}|pwk}{ }\textcolor{violet}{Premiere}\eventindex{Théâtre de la Porte Saint-Martin@\textbf{Théâtre de la Porte Saint-Martin}!Premiere von La femme nue@Premiere von La femme nue|pwkv}, das bereits 1908 am \textcolor{pink}{Théâtre de la Renaissance}\oindex{Théâtre de la Renaissance@\textbf{Théâtre de la Renaissance}, \emph{Theater}|pwk} seine \textcolor{violet}{Uraufführung}\eventindex{Théâtre de la Renaissance@\textbf{Théâtre de la Renaissance}!Uraufführung von La femme nue@Uraufführung von La femme nue|pwkv} erlebt hatte.
                  Die nächste Premiere von einem Schauspiel \textcolor{blue}{Batailles}\pwindex{Bataille, Henri 4.\,4.\,1872 Nîmes – 2.\,3.\,1922 Rueil-Malmaison@\textsc{Bataille, Henri} (4.\,4.\,1872 Nîmes – 2.\,3.\,1922 Rueil-Malmaison), \emph{Schriftsteller}|pwk} am \textcolor{pink}{Théâtre de la Porte
                     Saint-Martin}\oindex{Théâtre de la Porte Saint-Martin@\textbf{Théâtre de la Porte Saint-Martin}, \emph{Theater}|pwk}, die \textcolor{violet}{Uraufführung von \emph{\textcolor{green}{Les Flambeaux}\pwindex{Bataille, Henri 4.\,4.\,1872 Nîmes – 2.\,3.\,1922 Rueil-Malmaison@\textsc{Bataille, Henri} (4.\,4.\,1872 Nîmes – 2.\,3.\,1922 Rueil-Malmaison), \emph{Schriftsteller}!Flambeaux, pièce en trois actes@\strich\emph{Les Flambeaux, pièce en trois actes}|pwk}}}\eventindex{Théâtre de la Porte Saint-Martin@\textbf{Théâtre de la Porte Saint-Martin}!Uraufführung von Les Flambeaux@Uraufführung von Les Flambeaux|pwk}, gab es erst in der nächsten Spielzeit am 26. 11. 1912.}}}\label{K_L04009-2}.
               Lauter Lustspiele die leicht zu montieren {\pb}sind, gar nichts kosten. Der \label{K_L04009-3v}\edtext{\begin{otherlanguage}{french}Eschek\end{otherlanguage}}{\lemma{\textnormal{\emph{Eschek}}}\Cendnote{\textnormal{\begin{otherlanguage}{french}l’échec\end{otherlanguage}, französisch: Misserfolg; jiddisch heisik/hêsek: Schaden.}}}\label{K_L04009-3} des \label{K_L04009-4v}\edtext{\textcolor{green}{Chantekler}\pwindex{Rostand, Edmond 1.\,4.\,1868 Marseille – 2.\,12.\,1918 Paris@\textsc{Rostand, Edmond} (1.\,4.\,1868 Marseille – 2.\,12.\,1918 Paris), \emph{Schriftsteller}!Chantecler@\strich\emph{Chantecler}|pw}{}\ledrightnote{\textcolor{green}{Chantecler}}}{\lemma{\textnormal{\emph{Chantekler}}}\Cendnote{\textnormal{Die \textcolor{violet}{Uraufführung}\eventindex{Théâtre de la Porte Saint-Martin@\textbf{Théâtre de la Porte Saint-Martin}!Uraufführung von Chantecler@Uraufführung von Chantecler|pwkv} des neuen Stückes \emph{\textcolor{green}{Chantecler}\pwindex{Rostand, Edmond 1.\,4.\,1868 Marseille – 2.\,12.\,1918 Paris@\textsc{Rostand, Edmond} (1.\,4.\,1868 Marseille – 2.\,12.\,1918 Paris), \emph{Schriftsteller}!Chantecler@\strich\emph{Chantecler}|pwk}} des Erfolgsautors von \emph{\textcolor{green}{Cyrano de Bergerac}\pwindex{Rostand, Edmond 1.\,4.\,1868 Marseille – 2.\,12.\,1918 Paris@\textsc{Rostand, Edmond} (1.\,4.\,1868 Marseille – 2.\,12.\,1918 Paris), \emph{Schriftsteller}!Cyrano de Bergerac. Comédie héroïque en cinq actes en vers@\strich\emph{Cyrano de Bergerac. Comédie héroïque en cinq actes en vers}|pwk}}, \textcolor{blue}{Edmond Rostand}\pwindex{Rostand, Edmond 1.\,4.\,1868 Marseille – 2.\,12.\,1918 Paris@\textsc{Rostand, Edmond} (1.\,4.\,1868 Marseille – 2.\,12.\,1918 Paris), \emph{Schriftsteller}|pwk} im Vorjahr, war mit Spannung erwartet und mehrfach
                  verschoben worden. Die Inszenierung des \textcolor{green}{Dramas}\pwindex{Rostand, Edmond 1.\,4.\,1868 Marseille – 2.\,12.\,1918 Paris@\textsc{Rostand, Edmond} (1.\,4.\,1868 Marseille – 2.\,12.\,1918 Paris), \emph{Schriftsteller}!Chantecler@\strich\emph{Chantecler}|pwkv}, das über 70 Charaktere und noch mehr Kostüme
                  vorsah, war sehr aufwendig.}}}\label{K_L04009-4} hat \textcolor{blue}{Herz}\pwindex{Hertz, Henri 17.\,6.\,1875 Nogent-sur-Seine – 11.\,10.\,1966 16. arrondissement [Paris]@\textsc{Hertz, Henri} (17.\,6.\,1875 Nogent-sur-Seine – 11.\,10.\,1966 16. arrondissement [Paris]), \emph{Schriftsteller, Journalist, Theaterdirektor}|pw}{}\ledrightnote{\textcolor{blue}{Henri Hertz}} u \textcolor{blue}{Coquelin}\pwindex{Coquelin, Jean 1.\,12.\,1865 Paris – 1.\,10.\,1944@\textsc{Coquelin, Jean} (1.\,12.\,1865 Paris – 1.\,10.\,1944), \emph{Schauspieler}|pw}{}\ledrightnote{\textcolor{blue}{Jean Coquelin}} abgeschreckt – das
               grosse Schauspiel zu pflegen, u sie suchen ihr \textcolor{brown}{Théater}\orgindex{Théâtre de la Porte Saint-Martin@Théâtre de la Porte Saint-Martin|pwv}{}\ledrightnote{{$\rightarrow$}\emph{\textcolor{brown}{Théâtre de la Porte Saint-Martin}}} auf eine kleinere Basis zu stellen. Infolge dessen
               wurde ein grosser Teil des Personals als überflüssig entlassen. Sie müssten für den
                  {\pb}\textcolor{green}{Medardus}\pwindex{Schnitzler, Arthur 15. 5. 1862 Wien – 21. 10. 1931 ebd.@\textsc{Schnitzler, Arthur} (15. 5. 1862 Wien – 21. 10. 1931 ebd.), \emph{Schriftsteller, Mediziner}!junge Medardus. Dramatische Historie in einem Vorspiel und fünf Aufzügen@\strich\emph{Der junge Medardus. Dramatische Historie in einem Vorspiel und fünf Aufzügen}|pw}{}\ledrightnote{\textcolor{green}{Der junge Medardus. Dramatische Historie in einem Vorspiel und fünf Aufzügen}} (der jeden Falls erst Oktober 1912 hätte dranko{\geminationm}en können) Alles wieder
               anders einrichten. Trotzdem liessen mich die \textcolor{blue}{Direktoren}\pwindex{Hertz, Henri 17.\,6.\,1875 Nogent-sur-Seine – 11.\,10.\,1966 16. arrondissement [Paris]@\textsc{Hertz, Henri} (17.\,6.\,1875 Nogent-sur-Seine – 11.\,10.\,1966 16. arrondissement [Paris]), \emph{Schriftsteller, Journalist, Theaterdirektor}|pwv}\pwindex{Coquelin, Jean 1.\,12.\,1865 Paris – 1.\,10.\,1944@\textsc{Coquelin, Jean} (1.\,12.\,1865 Paris – 1.\,10.\,1944), \emph{Schauspieler}|pwv}{}\ledrightnote{{$\rightarrow$}\emph{\textcolor{blue}{Henri Hertz}}{\newline}{$\rightarrow$}\emph{\textcolor{blue}{Jean Coquelin}}} auf diesen Bescheid
               wochenlang warten. Weil wie mir \textcolor{blue}{Albert Clemen{\pb}ceau}\pwindex{Clemenceau, Albert 23.\,2.\,1861 Nantes – 23.\,7.\,1955 Sceaux@\textsc{Clemenceau, Albert} (23.\,2.\,1861 Nantes – 23.\,7.\,1955 Sceaux), \emph{Politiker, Jurist}|pw}{}\ledrightnote{\textcolor{blue}{Albert Clemenceau}} noch am Tag vor meiner Abreise sagte – die grösste Lust besteht \textcolor{green}{Medardus}\pwindex{Schnitzler, Arthur 15. 5. 1862 Wien – 21. 10. 1931 ebd.@\textsc{Schnitzler, Arthur} (15. 5. 1862 Wien – 21. 10. 1931 ebd.), \emph{Schriftsteller, Mediziner}!junge Medardus. Dramatische Historie in einem Vorspiel und fünf Aufzügen@\strich\emph{Der junge Medardus. Dramatische Historie in einem Vorspiel und fünf Aufzügen}|pw}{}\ledrightnote{\textcolor{green}{Der junge Medardus. Dramatische Historie in einem Vorspiel und fünf Aufzügen}} zu bringen. In der \textcolor{pink}{Auvergne}\oindex{Auvergne@\textbf{Auvergne}|pw}{}\ledrightnote{\textcolor{pink}{Auvergne}}
               wo beide \textcolor{blue}{Direktoren}\pwindex{Coquelin, Jean 1.\,12.\,1865 Paris – 1.\,10.\,1944@\textsc{Coquelin, Jean} (1.\,12.\,1865 Paris – 1.\,10.\,1944), \emph{Schauspieler}|pwv}\pwindex{Hertz, Henri 17.\,6.\,1875 Nogent-sur-Seine – 11.\,10.\,1966 16. arrondissement [Paris]@\textsc{Hertz, Henri} (17.\,6.\,1875 Nogent-sur-Seine – 11.\,10.\,1966 16. arrondissement [Paris]), \emph{Schriftsteller, Journalist, Theaterdirektor}|pwv}{}\ledrightnote{{$\rightarrow$}\emph{\textcolor{blue}{Jean Coquelin}}{\newline}{$\rightarrow$}\emph{\textcolor{blue}{Henri Hertz}}} ihr Ferien zubringen, scheint nun leider doch der mir mitgeteilte
               Entschluss gereift zu sein.\pend
           
\pstart
           {\pb}Es ist mir so leid Ihnen die Mühe des \textcolor{green}{Scenario}\pwindex{Schnitzler, Arthur 15. 5. 1862 Wien – 21. 10. 1931 ebd.@\textsc{Schnitzler, Arthur} (15. 5. 1862 Wien – 21. 10. 1931 ebd.), \emph{Schriftsteller, Mediziner}!junge Medardus. (Zur Übersetzung ins Französische)@\strich\emph{Der junge Medardus. (Zur Übersetzung ins Französische)}|pwv}{}\ledrightnote{{$\rightarrow$}\emph{\textcolor{green}{Der junge Medardus. (Zur Übersetzung ins Französische)}}} verursacht zu haben.
               Vielleicht ist dieselbe aber nicht ganz verloren. D\textsuperscript{r}{ }\textcolor{blue}{Frischauer}\pwindex{Frischauer, Berthold 9.\,9.\,1851 Brünn – 4.\,2.\,1924 Wien@\textsc{Frischauer, Berthold} (9.\,9.\,1851 Brünn – 4.\,2.\,1924 Wien), \emph{Journalist}|pw}{}\ledrightnote{\textcolor{blue}{Berthold Frischauer}} schreibt mir eben – ich möge meine
                  \introOben{}\textcolor{blue}{Herz}\pwindex{Hertz, Henri 17.\,6.\,1875 Nogent-sur-Seine – 11.\,10.\,1966 16. arrondissement [Paris]@\textsc{Hertz, Henri} (17.\,6.\,1875 Nogent-sur-Seine – 11.\,10.\,1966 16. arrondissement [Paris]), \emph{Schriftsteller, Journalist, Theaterdirektor}|pw}{}\ledrightnote{\textcolor{blue}{Henri Hertz}}\introOben{} gelieferte Übersetzung desselben zurückverlangen – da eine andere Kombination
               in \textcolor{pink}{Paris}\oindex{Paris@\textbf{Paris}, \emph{Hauptstadt}|pw}{}\ledrightnote{\textcolor{pink}{Paris}} möglich ist. Wenn Sie also Geduld {\pb}haben wollen – so könnten wir abwarten.\pend
           
\pstart
           Wegen des \textcolor{green}{weiten Landes}\pwindex{Schnitzler, Arthur 15. 5. 1862 Wien – 21. 10. 1931 ebd.@\textsc{Schnitzler, Arthur} (15. 5. 1862 Wien – 21. 10. 1931 ebd.), \emph{Schriftsteller, Mediziner}!weite Land. Tragikomödie in fünf Akten@\strich\emph{Das weite Land. Tragikomödie in fünf Akten}|pw}{}\ledrightnote{\textcolor{green}{Das weite Land. Tragikomödie in fünf Akten}} dencke ich an \textcolor{blue}{Gémier}\pwindex{Gémier, Firmin 21.\,2.\,1865 Aubervilliers – 26.\,11.\,1933 Paris@\textsc{Gémier, Firmin} (21.\,2.\,1865 Aubervilliers – 26.\,11.\,1933 Paris), \emph{Theaterleiter, Schauspieler, Drehbuchautor}|pw}{}\ledrightnote{\textcolor{blue}{Firmin Gémier}}. Ich werde wol im
                  Spätherbst wieder in \textcolor{pink}{Paris}\oindex{Paris@\textbf{Paris}, \emph{Hauptstadt}|pw}{}\ledrightnote{\textcolor{pink}{Paris}} sein –
               {\kaufmannsund} bin bereits (ohne das er weiss weshalb) in Relation mit ihm. Dafür muss ich nun
               aber noch Ihre besondere Autorisation haben –. Ende September also in
                  \textcolor{pink}{Wien}\oindex{Wien@\textbf{Wien}, \emph{Verwaltungsgebiet}|pw}{}\ledrightnote{\textcolor{pink}{Wien}}. Bitte verzeihen Sie – {\pb}\label{T_L04009-1v}\edtext{dass es nicht gelang –. Ihrer Frau \textcolor{blue}{Gemahlin}\pwindex{Schnitzler, Olga 17.\,1.\,1882 Wien – 13.\,1.\,1970 Lugano@\textsc{Schnitzler, Olga} (17.\,1.\,1882 Wien – 13.\,1.\,1970 Lugano), \emph{Schauspielerin, Sängerin}|pw}{}\ledrightnote{\textcolor{blue}{Olga Schnitzler}} u Ihnen herzlichste Grüsse. Ihre
               ergebene \spacefill\mbox{B. Zuckerkandl}}{\lemma{\textnormal{\emph{dass … Zuckerkandl}}}\Cendnote{\textnormal{Der Briefschluss befindet sich am unteren
                  Rand der fünften Seite.}}}\label{T_L04009-1}\pend
           \selectlanguage{ngerman}\endnumbering\briefempfaengerindex{Schnitzler, Arthur@\textsc{Schnitzler, Arthur}!zzzZuckerkandl, Berta@\emph{von Berta Zuckerkandl}!1911-07-251@{25. 7. {[}1911{]}}|)be}\mylabel{L04009h}  \newcommand{\dateiname}{L04009}\newcommand{\titel}{Berta Zuckerkandl an Arthur Schnitzler, 25. 7. [1911]}\newcommand{\editorInnen}{Herausgegeben von Jahnke, SelmaMüller, Martin Anton}\input{../tex-inputs/latex-korrekturansicht-abspann}
